% Source: sources/returns/S6_04_MAPS_MANIFOLDS/sections/manifolds.tex:715–775
\paragraph{Theorem D.1.12 (Proved: a canonical cyclic action and its quotient).}
There is an explicit free action $\sigma_d$ of $\mathbb Z/d$ on
$M_d$, defined by the affine and toric data below, such that every
admissible multiplicity-$(3,4)$ member with $|p|=d$ is
orientation-preservingly diffeomorphic to $M_d/\langle\sigma_d\rangle$.
All cyclic actions arising from these members are equivariantly
diffeomorphic to this action, allowing an automorphism of the
abstract deck group. Their quotients carry exactly the complex
structures and fixed-$u$ count in D.1.7.
\begin{proof}
Choose the specific original datum
$(\ell_0,\ell_1,\ell_2)=(0,-d,d)$ and scalar period $c=-i$.
It is admissible and has $p=d$. Its cover has lattice $\Lambda_d$
and matrices $A_1^{(d)},A_2^{(d)}$ of D.1.4. Use the respective
finite invariant translation vectors
\[
 v'_1=(-1,-2d,4d,0),\qquad v'_2=(1,3d,-3d,0),
\]
divided by three and four in the affine finite actions, and use
the cusp period lattice $\langle e_1,d e_2\rangle$. The invariant
meridian value is $p'=1$. This specifies a single smooth model,
independently of any input twist ambiguity; D.1.11 identifies its
underlying manifold with $M_d$.

On a regular torus define $\sigma_d$ by translation through the
original $\widehat\gamma$, whose coordinate in $\Lambda_d$ is
$(1/d,0,0,0)$. The matrices $A_j^{(d)}$ change this vector by an
integral vector, so this translation is globally defined on the
regular family and commutes with the finite quotient actions.
It has order $d$. On a finite reduced fibre a fixed point of its
$a$th power would require, in the invariant $\gamma'$ coordinate,
$a/d\equiv\pm k/m_j\pmod1$ for some $k$. Since $\gcd(d,m_j)=1$,
this forces $d\mid a$. On the cusp it is the affine period
translation by $-e_2$, including the prescribed holomorphic unit,
and cycles the $d$ components. No nonzero power below $d$ fixes a
stratum: after subtracting a covering period it would translate a
finite triangle, edge or vertex to itself by a nonzero vector.
Thus it extends over the entire toric tube and the action is free.
The quotient restores the original lattice, twists $(-d,d)$ and
period $c=-i$, so is exactly that fixed representative.

D.1.7 gives an oriented diffeomorphism from any other constructed
quotient with the same $d$ to this one. Lift it to universal covers.
Uniqueness of path lifting makes the lift intertwine deck
transformations by the induced automorphism of $\pi_1=\mathbb Z/d$.
This proves the claimed equivariant classification within the
entire specified construction by a direct proof, without assuming
a classification of arbitrary free actions on $M_d$. Different
$d$ are distinguished downstairs by fundamental group and upstairs
by Euler characteristic.

The fixed-$u$ count is two for the present multiplicities. In the
coordinates of Viaclovsky's higher-multiplicity sphere family,
the checked dictionary K1 and R4 instead give $2k\ell$ at fixed
$u$ for the pair $(3k,4\ell)$. K12 supplies tangent cohomology
$(1,1,1,0)$ and Kodaira dimension determined by the sign of
$k\ell-3k-4\ell$ for those sphere members; at $(k,\ell)=(1,1)$
this specializes to the count two and negative Kodaira dimension.
These sphere results concern $d=1$ and do not assert tangent
cohomology for the torsion covers $M_d$ with $d>1$.
\end{proof}
